\documentclass[10pt,letterpaper]{article}
\usepackage[margin=1in]{geometry}
\usepackage{amsmath,amssymb}
\usepackage{booktabs}
\usepackage{graphicx}
\usepackage{hyperref}
\usepackage{microtype}

\title{FastVLM as a Robot Vision--Language--Action Backbone:\\ An MLX-Native On-Device Study}
\author{Syun Robotics Lab}
\date{Draft --- 2026-11-01}

\begin{document}
\maketitle

\begin{abstract}
We study whether FastVLM's efficient high-resolution visual encoding remains useful after
adaptation to robotic action prediction. We implement a fully MLX-native FastVLA system on an
Apple M4 Mac and train an action-chunking policy over FastVLM representations, with frozen and
language-LoRA variants. Controlled comparisons cover state-only, FastViTHD-only, global-only,
FastVLM frozen, FastVLM LoRA, and a SmolVLM2-500M backbone on episode-disjoint Meta-World data.
Evaluation measures native task success, manipulation error, action stability, instruction
counterfactuals, latency, memory, and training efficiency. All reported experimental values and
statistical claims are generated from the files under \texttt{results/confirmatory}; this draft
must not be treated as complete until those placeholders are replaced by measured values.
\end{abstract}

\section{Introduction}
Vision-language-action (VLA) policies combine pretrained visual-language representations with
robot action decoding \cite{brohan2023rt2,kim2024openvla,oxe2023}. FastVLM targets efficient
high-resolution vision-language inference by reducing visual tokens with FastViTHD
\cite{fastvlm2025}. It is unclear whether that efficiency translates to manipulation, where
small spatial errors and temporal action consistency matter more than text generation.

We ask three questions. (1) Does a frozen FastVLM representation support accurate closed-loop
control? (2) Does real MLX LoRA adaptation improve a controlled action head? (3) How do task
success, latency, and memory compare with state-only, visual-token, and lightweight-VLM
alternatives on the same demonstrations and simulator?

\section{Related work}
RT-2 and OpenVLA demonstrate large pretrained vision-language models as action policies
\cite{brohan2023rt2,kim2024openvla}. ACT and diffusion policies show the importance of action
chunking and temporal modeling \cite{zhao2023act,chi2023diffusion}. FastVLM instead focuses on
reducing visual encoding cost while retaining high input resolution \cite{fastvlm2025}. LoRA
provides a parameter-efficient adaptation mechanism \cite{hu2021lora}. Our study is deliberately
small and controlled: it isolates the backbone and visual-token path while keeping the action
head, data, action horizon, and simulator evaluation protocol fixed.

\section{Method}
\subsection{MLX-native multimodal policy}
The FastVLM backbone performs its own 1024-pixel FastViTHD preprocessing, visual encoding,
two-layer multimodal projection, and Qwen2 language processing in MLX. For an observation image
$I_t$, instruction $L$, and proprioceptive state $s_t$, the model produces a final language hidden
state $h_t$ and 256 projected FastViTHD tokens $V_t$. We average-pool $V_t$ from a $16\times16$
grid to a $4\times4$ spatial summary, project each spatial token, and concatenate it with $h_t$
and a projected normalized state. A two-layer MLP predicts an $H\times A$ action chunk.
The action dimension, state dimension, camera list, names, and horizon are configuration fields;
none are Meta-World constants in the core model.

\subsection{Training variants}
\paragraph{Frozen FastVLM.} The image tower, projector, and language model are frozen. Projected
visual tokens and the final hidden state are cached from the actual Apple checkpoint. Only the
action head is optimized with AdamW and masked action-chunk MSE.

\paragraph{FastVLM LoRA.} Upstream MLX \texttt{LoRaLayer} modules replace Qwen2 attention and MLP
projections. FastViTHD and the multimodal projector remain frozen, while the LoRA $A/B$ matrices
and action head are optimized jointly with \texttt{mlx.value\_and\_grad}. We verify trainable
parameter counts, nonzero gradients, and parameter changes before treating a run as LoRA.

\subsection{Data and simulation}
We use \texttt{lerobot/metaworld\_mt50} (Apache-2.0) and Meta-World v3 \cite{finn2019metaworld}.
Episodes are split 35/7/8 per selected task with a fixed seed. Action targets are eight consecutive
80-Hz actions with episode-boundary masks. State and action normalization use training episodes
only. Closed-loop evaluation uses the benchmark's native success flag and a fresh seed range.

\section{Experiments}
% Generated from measured result files only. Replace this block via scripts/build_paper.py
\subsection{Closed-loop success}
\textbf{PENDING: no confirmatory multi-seed result has been inserted into this draft yet.}

\subsection{Efficiency}
\textbf{PENDING: no confirmatory benchmark table has been inserted into this draft yet.}


\section{Discussion and limitations}
This study is simulation-only and does not establish physical-robot performance. The three tasks
do not cover all manipulation regimes, and the selected pretrained histories differ across
backbones. FastViTHD and the projector remain frozen in cached-feature and LoRA conditions. We use
three training seeds and pooled held-out episodes with Wilson intervals, but a negative or
inconclusive result remains a valid scientific outcome. We do not claim State-of-the-Art
performance.

\section{Reproducibility, ethics, and licensing}
The repository records split manifests, configurations, Git commits, W\&B run identifiers,
checkpoint metadata, and runtime versions. Apple FastVLM weights are licensed under
\texttt{apple-amlr} for non-commercial research. This work is an unofficial derivative and has no
Apple endorsement. Dataset bytes are referenced under their original license rather than
redistributed. No API credentials are committed.

\bibliographystyle{plain}
\bibliography{references}
\end{document}
